\documentclass[11pt]{article}
\usepackage[margin=1.1in]{geometry}
\usepackage{amsmath, amssymb, amsthm, mathtools}
\usepackage{enumitem}
\usepackage[colorlinks=true, linkcolor=blue, citecolor=blue]{hyperref}

% Long tags for the class conditions: a plain itemize sets \item[...] labels
% right-aligned in a fixed-width box, so anything longer than \labelwidth runs
% off the left margin. `style=unboxed' lets the tag flow into the body instead.
\newlist{conditions}{description}{1}
\setlist[conditions]{style=unboxed, leftmargin=0pt, labelsep=0.35em,
                     itemsep=5pt, parsep=2pt, font=\normalfont}

\newtheorem{theorem}{Theorem}
\newtheorem{lemma}{Lemma}
\newtheorem{proposition}{Proposition}
\newtheorem{corollary}{Corollary}
\newtheorem{assumption}{Assumption}
\theoremstyle{definition}
\newtheorem{definition}{Definition}
\newtheorem{remark}{Remark}
\newtheorem{example}{Example}

\newcommand{\R}{\mathbb{R}}
\newcommand{\E}{\mathbb{E}}
\newcommand{\cA}{\mathcal{A}}
\newcommand{\cB}{\mathcal{B}}
\newcommand{\cP}{\mathcal{P}}
\newcommand{\cN}{\mathcal{N}}
\newcommand{\ip}[2]{\left\langle #1, #2 \right\rangle}
\newcommand{\KL}{\mathrm{KL}}
\DeclareMathOperator*{\argmax}{arg\,max}
\DeclareMathOperator{\supp}{supp}
\DeclareMathOperator{\conv}{conv}

\title{A Game Class for Convergence of Magnetic Mirror Descent\\
on Gaussian-Mixture Policies}
\author{\texttt{cont\_tinkering} / \texttt{theory2}}
\date{July 2026}

\begin{document}
\maketitle

\begin{abstract}
We define a class of continuous-action two-player zero-sum games on which
Magnetic Mirror Descent (MMD) with a $K$-component Gaussian-mixture policy is
expected to converge to a Nash equilibrium, under the standing assumptions that
the equilibrium has finite support of size $M$ and that $K \ge M$. The class is
built from a single observation: with a mixture policy, the categorical head
plays a finite matrix game \emph{exactly}, and tabular MMD solves finite matrix
games globally; therefore every obstruction to convergence lives in the
\emph{mean} dynamics, which are natural-gradient ascent on a smoothed effective
landscape. The class accordingly imposes one condition --- a uniform
one-point-attraction (star-concavity) condition on the family of smoothed,
magnet-regularised effective landscapes, uniformly over opponent policies and
over the admissible range of component scales --- together with capacity,
identifiability and step-size conditions. This one condition subsumes the three
empirically separate failure edges of the previous theory (decoy capture,
coupling ejection, dead-zone freeze), is stated in a way that is dimension-free
and asymmetric in the two players, and is sharp in the sense that its qualitative
failure produces an exponentially stable non-Nash attractor with an open basin.
The conditions are imposed on scales bounded below by the implementation's
variance floor, which is what allows the payoff itself to be merely bounded and
measurable and what makes the class tolerant of landscape structure finer than
the floor; exactness is recovered separately, as a corollary in which the floor
is released and the scale dynamics drive the representation bias to zero at rate
$\Theta(1/t)$ --- a prediction confirmed numerically here, together with a closed
form for how the Gaussian magnet rescales it. Section~\ref{sec:suff} gives the argument for why the class suffices,
in the form of a three-block input-to-state-stability decomposition intended as
the skeleton of a future proof; Section~\ref{sec:check} gives checkable
sufficient conditions (finite-rank couplings, Gaussian-bump landscapes, separable
multidimensional games).
\end{abstract}

\tableofcontents

\section{Setting}

\subsection{Game}

Player $0$ maximises and player $1$ minimises a payoff
$u : \cA \times \cB \to \R$, where $\cA \subset \R^{n_0}$ and
$\cB \subset \R^{n_1}$ are compact convex bodies and $u$ is \emph{bounded and
measurable} (Remark~\ref{rem:reg} explains why nothing stronger is needed).
\emph{No symmetry between the players is assumed}: the action sets,
dimensions, landscapes, support sizes, component counts and all hyperparameters
may differ. Mixed strategies are $\mu \in \cP(\cA)$, $\nu \in \cP(\cB)$ and
$U(\mu,\nu) = \E_{\mu \otimes \nu} u$.

The \emph{effective landscapes} are
\begin{equation}\label{eq:eff}
  Q_\nu(a) \coloneqq \int u(a,b)\,d\nu(b), \qquad
  W_\mu(b) \coloneqq -\int u(a,b)\,d\mu(a),
\end{equation}
so both players \emph{maximise}: player $0$ maximises $\E_\mu Q_\nu$, player $1$
maximises $\E_\nu W_\mu$. For a function $\Lambda$ on $\R^n$ and a scale vector
$s \in \R^n_{>0}$ we write $\Lambda_s \coloneqq \Lambda * \cN(0, \operatorname{diag}(s^2))$.

We fix a mixed Nash equilibrium $(\mu^\star,\nu^\star)$ and assume throughout ---
this is a standing assumption, not something we derive --- that both equilibrium
measures are \emph{purely atomic with finite support}:
\begin{equation}\label{eq:finite}
  \mu^\star = \sum_{j=1}^{M_0} w^\star_j\, \delta_{p_j},
  \qquad
  \nu^\star = \sum_{l=1}^{M_1} v^\star_l\, \delta_{r_l},
  \qquad
  w^\star \in \Delta_{M_0},\ v^\star \in \Delta_{M_1},\ w^\star, v^\star > 0 .
\end{equation}
\textbf{The points $p_1,\dots,p_{M_0} \in \cA$ and $r_1,\dots,r_{M_1} \in \cB$
are called the \emph{atoms}} of the equilibrium (the term is the measure-theoretic
one: an atom of a measure is a point carrying strictly positive mass). They are
distinct, $M_i$ denotes their number, and
\[
  \supp\mu^\star = \{\,p_j : j = 1,\dots,M_0\,\}, \qquad
  \supp\nu^\star = \{\,r_l : l = 1,\dots,M_1\,\}
\]
--- throughout, $j$ indexes player $0$'s atoms and $l$ indexes player $1$'s, and
both are bound indices ranging over the whole set. The two atom sets are
unrelated: they live in different spaces and $M_0 \ne M_1$ is allowed. Finding the atoms --- their
locations, via the component means --- is what the mean dynamics must do; finding
their weights $w^\star, v^\star$ is what the categorical heads must do.

\begin{remark}[The atoms are automatically maxima]\label{rem:support}
By the support characterisation of a best response,
$\supp\mu^\star \subseteq \argmax_{\cA} Q_{\nu^\star}$ and
$\supp\nu^\star \subseteq \argmax_{\cB} W_{\mu^\star}$. So ``the equilibrium
atoms sit at maxima of the effective landscape'' is a theorem, not a hypothesis.
Everything the class must assume about the atoms is (i) nondegeneracy of those
maxima and (ii) that the landscape has \emph{no other} local maxima --- see
Definition~\ref{def:class}(C) and Proposition~\ref{prop:sharp}.
\end{remark}

\subsection{Policy and update}

Player $i$ uses a $K_i$-component diagonal-Gaussian mixture
\begin{equation}
  \pi_{\theta_i} = \sum_{k=1}^{K_i} w^{(i)}_k\,
  \cN\!\big(m^{(i)}_k, \operatorname{diag}(\sigma^{(i)}_k)^2\big),
  \qquad w^{(i)} = \mathrm{softmax}(\ell^{(i)}),
\end{equation}
with each $\sigma^{(i)}_{k,d}$ clipped to a range
$\Sigma_i \coloneqq [s^{\min}_i, s^{\max}_i]$ (both bounds exist in the
implementation: \texttt{idealized\_mmd.py} clips $\log\sigma$ to
$[\log 10^{-3}, \log 1]$). The update, per player, per step, with step size
$\eta_i$, categorical magnet temperature $\tau_i^{\mathrm{cat}}$, Gaussian magnet
temperature $\tau_i^{\mathrm{g}} \ge 0$, and magnet snapshot $\bar\theta_i$
refreshed every $T$ steps, is
\begin{align}
  \ell^{(i),+} &= \frac{\eta_i q^{(i)} + \eta_i \tau^{\mathrm{cat}}_i
    \log \bar w^{(i)} + \log w^{(i)}}{1 + \eta_i \tau^{\mathrm{cat}}_i},
  \label{eq:cat}\\
  m^{(i),+}_k &= m^{(i)}_k + \eta_i (\sigma^{(i)}_k)^2 \odot
    \nabla_{m_k} \Phi_i, \qquad
  \rho^{(i),+}_k = \rho^{(i)}_k + \tfrac{\eta_i}{2} \nabla_{\rho_k} \Phi_i,
  \label{eq:gauss}
\end{align}
where $\rho = \log\sigma$ and
$\Phi_i = (\pm)U - \tau^{\mathrm{g}}_i \KL(\pi_{\theta_i} \| \pi_{\bar\theta_i})$.

The floor $s^{\min}_i > 0$ is \emph{strict} and is used throughout. It is not a
numerical detail: it is the hypothesis that makes the class both broad and
checkable, as the next remark and Section~\ref{sec:floor} explain.

\begin{remark}[Why bounded measurable suffices]\label{rem:reg}
Every landscape the algorithm evaluates is smoothed at a scale bounded below by
$s^{\min}_i > 0$. For bounded measurable $u$ and any $s \ge s^{\min}$,
differentiating under the integral sign onto the Gaussian kernel gives
\begin{equation}\label{eq:derivbound}
  \big\|\nabla^k (Q_\nu)_s\big\|_\infty \;\le\;
  C_k\, \|u\|_\infty \,/\, (s^{\min})^{k} \;<\; \infty
  \qquad \text{for every } k,
\end{equation}
uniformly in $\nu$ and in $s \in \Sigma_0$. So $(Q_\nu)_s$ is $C^\infty$ with
uniformly bounded derivatives no matter how irregular $u$ is, and every
hypothesis of Definition~\ref{def:class} --- all of which are stated on smoothed
objects --- has finite constants. This admits payoffs that are not $C^2$ and not
even continuous: $|a-b|$, $\min(a,b)$, hard-threshold Blotto payoffs, indicators.
Smoothness of $u$ is needed at exactly one point in the development, the
\emph{exponent} of the representation bias, and only near the equilibrium atoms;
it is isolated as condition (B) below. Note the price of dropping the floor:
\eqref{eq:derivbound} degrades as $s^{\min} \to 0$, and the hypotheses would have
to be imposed on the raw payoff.
\end{remark}

\subsection{The structural fact the whole class is built on}
\label{sec:structure}

\begin{lemma}[Exact decomposition of the update]\label{lem:decomp}
For any bounded measurable $u$, define the \emph{component matrix}
\begin{equation}\label{eq:Abar}
  \bar A_{kl}(m,\sigma) \coloneqq
  \E_{a \sim \cN(m^{(0)}_k, \sigma^{(0)2}_k),\; b \sim \cN(m^{(1)}_l, \sigma^{(1)2}_l)}
  \big[u(a,b)\big], \qquad
  \bar A \in \R^{K_0 \times K_1}.
\end{equation}
Then
\begin{equation}
  U(\pi_{\theta_0}, \pi_{\theta_1}) = (w^{(0)})^\top \bar A\, w^{(1)}
  = \sum_k w^{(0)}_k\, (Q_{\pi_{\theta_1}})_{\sigma_k}\!\big(m^{(0)}_k\big),
\end{equation}
and consequently
\begin{enumerate}
\item[(i)] $q^{(0)} = \bar A w^{(1)}$ and $q^{(1)} = -\bar A^\top w^{(0)}$, so
  \eqref{eq:cat} is \emph{verbatim} tabular MMD on the finite zero-sum matrix
  game $\bar A$;
\item[(ii)] the mean update \eqref{eq:gauss} is
  $m^{(0),+}_k = m^{(0)}_k + \eta_0\, \sigma_k^2 w^{(0)}_k
   \nabla (Q_{\pi_{\theta_1}})_{\sigma_k}(m^{(0)}_k) + (\text{magnet term})$,
  i.e.\ preconditioned gradient ascent on the smoothed effective landscape at the
  component's own scale.
\end{enumerate}
\end{lemma}

\begin{proof}
Immediate from Fubini and the definition of the mixture; (i) is the observation
that $Q$ is linear in $\nu$ and the component expectations are exactly the matrix
entries; (ii) is $\nabla_{m_k} U = w_k \nabla (Q_\nu)_{\sigma_k}(m_k)$ together
with the per-component inverse Fisher factor $\sigma_k^2$.
\end{proof}

Lemma~\ref{lem:decomp}(i) is the decisive one: \textbf{the categorical heads
solve a finite matrix game, and tabular MMD on a finite zero-sum matrix game
converges last-iterate from any initialisation} (the lifted game is bilinear,
hence monotone; Sokota et al.\ 2023). There is no basin structure and no
initialisation dependence in that block. Therefore:

\begin{center}
\emph{Every obstruction to global convergence of mixture-MMD lives in the mean
dynamics of Lemma~\ref{lem:decomp}(ii).}
\end{center}

The class below consists of exactly the hypotheses needed to make the mean
dynamics find the atoms, plus the standard bookkeeping (capacity,
identifiability, step size).

\begin{remark}[The $w_k$ factor is an implementation artifact]\label{rem:wk}
The factor $w_k$ in Lemma~\ref{lem:decomp}(ii) is \emph{not} part of the natural
gradient of the mixture policy. The Fisher metric of the hierarchical law
$p(k,a) = w_k \cN(a \mid m_k, \sigma_k^2)$ has $m_k$-block $w_k/\sigma_k^2$, so
the exact natural gradient is $\sigma_k^2 \nabla (Q_\nu)_{\sigma_k}(m_k)$, with no
$w_k$. The implemented preconditioner $\sigma_k^2$ is the inverse Fisher metric
of the \emph{conditional} law only. The discrepancy is the entire
weight-starvation freeze: a component losing weight also loses mobility and
freezes at a non-atom. We carry hypothesis~(W) of Definition~\ref{def:class} to handle it,
and note that the hypothesis is vacuous if the preconditioner is corrected to
$\sigma_k^2/w_k$.
\end{remark}

\section{The class}

\subsection{Auxiliary objects}

\paragraph{Reachable policy set.} Let $\cN_1 \subseteq \cP(\cB)$ be a set of
opponent policies containing $\nu^\star$, the initial policy, and every policy
the run visits. Any forward-invariant superset works; the safe default is
$\cN_1 = \cP(\cB)$, and under a finite-rank coupling the quantifier
``for all $\nu \in \cN_1$'' collapses to a quantifier over a compact convex
subset of $\R^m$ (Section~\ref{sec:check}). Define $\cN_0 \subseteq \cP(\cA)$
symmetrically.

\paragraph{Regularised landscape.} For a magnet centre $\bar a$ and Gaussian
magnet temperature $\tau^{\mathrm g}_0$, put
\begin{equation}\label{eq:reg}
  Q^{\mathrm{reg}}_{\nu, s, \bar a}(a) \coloneqq (Q_\nu)_s(a)
  \;-\; \frac{\tau^{\mathrm g}_0}{2}\,\big\|a - \bar a\big\|^2_{\bar\sigma^{-2}} ,
\end{equation}
the quadratic term being the mean-dependent part of
$\KL\big(\cN(a,\sigma^2)\,\|\,\cN(\bar a,\bar\sigma^2)\big)$. With
$\tau^{\mathrm g}_0 = 0$ this is just the smoothed effective landscape;
allowing $\tau^{\mathrm g}_0 > 0$ is what lets the class cover games whose
landscape supplies no curvature at all (Remark~\ref{rem:flat}).

\paragraph{Star-concavity modulus.} A continuous $\rho : [0,\infty) \to
[0,\infty)$ with $\rho(0) = 0$ and $\rho(r) > 0$ for $r > 0$.

\subsection{Definition}

\begin{definition}[Uniformly attracting game, $\mathfrak{G}(M,\Sigma,\cN,\rho)$]
\label{def:class}
The game $u$ belongs to the class if the following hold. We state them for
player $0$; player $1$'s conditions are the same statements applied to $W_\mu$,
$\cB$, $M_1$, $\Sigma_1$, $\cN_0$, with its own constants.

\begin{conditions}
\item[\textbf{(C) Uniform attracting cell structure.}] There are open sets
  $\cA_1,\dots,\cA_{M_0} \subset \cA$ with disjoint closures whose union has full
  measure in $\cA$, and \emph{branch maps}
  $\xi_j : \cN_1 \times \Sigma_0 \times \cA_j \to \cA_j$,
  $(\nu, s, \bar a) \mapsto \xi_j(\nu,s,\bar a)$, such that for every
  $\nu \in \cN_1$, every $s \in \Sigma_0^{n_0}$, every magnet centre
  $\bar a \in \cA_j$ and every $a \in \cA_j$,
  \begin{equation}\label{eq:star}
    \big\langle \nabla Q^{\mathrm{reg}}_{\nu,s,\bar a}(a),\;
      \xi_j - a \big\rangle \;\ge\; \rho\big(\|\xi_j - a\|\big)\,\|\xi_j - a\| .
  \end{equation}
  (Star-concavity, a.k.a.\ quasar-concavity, toward the branch point, uniformly
  in the opponent, the scale and the magnet.) When $\tau^{\mathrm g}_0 = 0$ the
  dependence on $\bar a$ is absent.
\item[\textbf{(D) Branch regularity and small gain.}] Each $\xi_j$ is Lipschitz
  in $(\nu, s, \bar a)$, with $\nu$ measured in a metric metrising weak
  convergence on $\cP(\cB)$ --- i.e.\ the weak-$*$ topology
  $\nu_k \to \nu \iff \int f \, d\nu_k \to \int f \, d\nu$ for all
  $f \in C(\cB)$, which is metrizable because $\cB$ is compact metric (any
  Lévy--Prokhorov or Wasserstein metric will do; the condition does not depend on
  the choice) --- and the round-trip gain of the branch maps is a contraction:
  \begin{equation}\label{eq:smallgain}
    g_0 \cdot g_1 < 1, \qquad
    g_0 \coloneqq \sup_j \operatorname{Lip}_\nu \xi_j, \quad
    g_1 \coloneqq \sup_l \operatorname{Lip}_\mu \zeta_l ,
  \end{equation}
  where $\zeta_l$ are player $1$'s branch maps and $\operatorname{Lip}_\nu$ is
  taken with respect to the displacement of the opponent's component means.
\item[\textbf{(M) Modal match.}] $M_0 = |\supp \mu^\star|$ and, after
  relabelling,
  \begin{equation}
    \xi_j\big(\nu^\star, s^{\min}_0, \cdot\big) = p_j + O\big((s^{\min}_0)^2\big)
    \qquad \text{for each } j :
  \end{equation}
  the branch endpoints at equilibrium and at the \emph{floor} scale are the
  equilibrium atoms, up to the smoothing displacement. (Anchoring at the floor
  rather than at scale $0$ is forced by the fact that the run never reaches
  scale $0$; the displacement is $O(s^2)$ under (B) and $O(s)$ without it.)
\item[\textbf{(B) Atom regularity} --- \emph{for the bias exponent only}.]
  $Q_{\nu^\star}$ is $C^2$ in a neighbourhood of each $p_j$, with
  $\nabla^2 Q_{\nu^\star}(p_j) \prec 0$. This condition is used \emph{nowhere}
  except to give the representation bias the exponent $2$ in
  Theorem~\ref{thm:main}; dropping it costs an order in the floor and nothing
  else (Remark~\ref{rem:biasexp}).
\item[\textbf{(I) Nondegeneracy.}] The $M_0 \times M_1$ matrix game
  $A_{jl} = u(p_j, r_l)$ --- the payoff table between the two players' atoms ---
  has a \emph{nondegenerate} equilibrium: equivalently, its Nash equilibrium is
  unique. Used only for well-posedness and to name the limit; see
  Remark~\ref{rem:nondeg}, which shows the conclusion survives its failure with
  ``a Nash of $A$'' in place of $(w^\star, v^\star)$.
\item[\textbf{(K) Capacity and covering.}] $K_i \ge M_i$, and the initial
  component means of player $i$ meet every cell: each $\cA_j$ contains at least
  one $m^{(0)}_k(0)$, each $\cB_l$ at least one $m^{(1)}_k(0)$. Initial scales lie
  in $\Sigma_i$.
\item[\textbf{(S) Algorithmic caps.}] $\tau^{\mathrm{cat}}_i > 0$; the step sizes
  $\eta_i$ are below the threshold set by the Lipschitz constant of the
  component-matrix game and of $\nabla Q^{\mathrm{reg}}$ on the relevant sets.
\item[\textbf{(W) Weight non-degeneracy.}] Along the trajectory,
  at least one component in each cell keeps weight $\ge w_{\min} > 0$. (Vacuous
  under the corrected preconditioner of Remark~\ref{rem:wk}; otherwise implied by
  (I) plus a categorical entropy floor $\tau^{\mathrm{ent}} > 0$, which bounds
  $w$ below by $\propto e^{-\mathrm{range}(q)/\tau^{\mathrm{ent}}}$.)
\end{conditions}
\end{definition}

\begin{lemma}[Resolution at the floor; a design rule for $s^{\min}$]
\label{lem:resolution}
Definition~\ref{def:class} forces the equilibrium atoms to be resolvable at the
bottom of the scale range: (C) provides $M_0$ cells with distinct branch points
at every $s \in \Sigma_0$, in particular at $s^{\min}_0$, and (M) identifies
those branch points with the atoms there. Hence
$(Q_{\nu^\star})_{s^{\min}_0}$ has $M_0$ distinct local maxima. For a landscape
of equal Gaussian bumps of width $\bar w_0$ this is the explicit inequality
\begin{equation}\label{eq:resolution}
  \sqrt{\bar w_0^2 + (s^{\min}_0)^2} \;<\; \tfrac12
  \min_{j \ne j'} \|p_j - p_{j'}\| .
\end{equation}
\end{lemma}

\begin{proof}
The first statement is the implication just given. For \eqref{eq:resolution}:
a bump of width $\bar w$ smoothed at scale $s$ has width
$W = \sqrt{\bar w^2 + s^2}$, and the sum of two equal Gaussians of width $W$ at
separation $d$ is bimodal if and only if $d > 2W$ --- the classical
separation threshold for an equal-weight, equal-variance Gaussian mixture.
Applying it to the closest pair of atoms gives \eqref{eq:resolution}.
\end{proof}

\begin{remark}[How to read Lemma~\ref{lem:resolution}]\label{rem:resolution}
It is stated separately not because it is an independent hypothesis --- it is
not, it is a consequence of (C) and (M) --- but because of its \emph{role}. Every
condition in Definition~\ref{def:class} constrains the game given the scale
range; read in the other direction, \eqref{eq:resolution} constrains a
\emph{hyperparameter} given the game: \emph{do not choose a floor larger than
half the atom spacing}. Definition~\ref{def:class} encodes that silently, and it
is the only part of the class with a one-line closed-form check.

Note also which term binds. For the codebase's two-peak game
($\bar w_0 = 0.1$, $s^{\min}_0 = 10^{-3}$, $d = 1$) the left side is
$\sqrt{10^{-2} + 10^{-6}} = 0.1000005$ against $d/2 = 0.5$: the floor's
contribution $10^{-6}$ is negligible beside the bump width $10^{-2}$. So in
practice \eqref{eq:resolution} is a statement about the \emph{landscape} --- the
peaks must be narrower than half their spacing, i.e.\ be visibly separate peaks
at all --- and holds with a factor $5$ to spare. The floor becomes the binding
term only when $s^{\min}$ is pushed up to the scale of the atom separation,
which is what an annealing schedule that never anneals down would do.
\end{remark}

\subsection{Reading the definition}

Condition (C) is the whole content; everything else is bookkeeping. Four
comments on what it does and does not say.

\begin{remark}[The target of (C) is an atom, not an equilibrium]
\label{rem:atomnotnash}
Since $M_0 > 1$, no single action is a Nash equilibrium, and \eqref{eq:star}
never claims one is: its target $\xi_j$ is \emph{one atom} --- one of the $M_0$
points the equilibrium mixes over --- not the equilibrium. The problem the class
addresses splits in two, and (C) is responsible for only the first:
\begin{center}
\begin{tabular}{@{}l p{0.36\textwidth} p{0.34\textwidth}@{}}
\hline
& \emph{where are the atoms?} & \emph{how much mass on each?}\\
\hline
solved by & the Gaussian heads & the categorical heads\\
mechanism & $M_i$ parallel single-target problems, one per cell, each governed
  by \eqref{eq:star} & one finite matrix game $\bar A$, solved globally by
  tabular MMD (Lemma~\ref{lem:decomp})\\
needs & condition (C) & condition (I)\\
\hline
\end{tabular}
\end{center}
This is why (C) is imposed \emph{cellwise on a partition} rather than globally:
the landscape has $M_0$ maxima, so no global star-concavity can hold, whereas on
each cell there is a single target and the standard one-point machinery applies
verbatim. The mixing is not a geometric phenomenon at all --- it is a property of
the weights, and the weights solve an ordinary matrix game.

The two halves meet at equilibrium in the following way. By
Remark~\ref{rem:support} every atom attains $\max Q_{\nu^\star}$, so the $M_0$
targets are \emph{tied}:
$Q_{\nu^\star}(\xi_1) = \dots = Q_{\nu^\star}(\xi_{M_0}) = \max Q_{\nu^\star}$.
Off equilibrium they are not tied; it is the \emph{opponent's} weights that tilt
the landscape until they are --- the classical indifference condition of a mixed
equilibrium. So each $\xi_j$ is a global maximiser at equilibrium, of a landscape
the opponent has arranged to be flat-topped across all $M_0$ cells.

This also locates \eqref{eq:star} against the literature it is borrowed from. As
an inequality it is standard: \emph{quasar-convexity} (equivalently weak
quasi-convexity) in nonconvex optimization, and \emph{variational stability} in
learning-in-games. The optimization line is single-target, static and
game-free --- we borrow the inequality and its Lyapunov argument, and apply them
$M_i$ times to a landscape that moves. The game-theoretic line \emph{does} cover
mixed equilibria, but only because there the strategy set is convex, so
$\mu^\star$ is a single point of it; that is the measure-space picture, in which
a one-point condition is immediately available and no partition is needed. The
parametric update does not live in that space --- the mixture manifold is
nonconvex and a component mean is a point of $\cA$, where no equilibrium
lives --- and the cellwise form of (C) is exactly the repair: descend to the level
at which the target \emph{is} single-valued, and recover the mixing from the
categorical head.
\end{remark}

\begin{remark}[What (I) is and is not for]\label{rem:nondeg}
Two clauses one might expect in (I) are free rather than assumed. That
$(w^\star, v^\star)$ is a Nash of $A$ is automatic --- $(\mu^\star,\nu^\star)$ is a
Nash of the continuous game, so in particular no reshuffling of mass among the
atoms helps --- and that it is fully mixed is already declared in
\eqref{eq:finite}, the atoms being the support. The whole content of (I) is
therefore nondegeneracy, and it is worth separating what that buys.

\emph{It is not needed for the conclusion.} Zero-sum equilibria are
\emph{interchangeable}: the optimal-strategy sets are convex, any pair of
optimal strategies is a Nash, and all have the same value. So converging to
\emph{some} equilibrium of $A$ is, in exploitability, exactly as good as
converging to a designated one; and Lemma~\ref{lem:weight} needs only
monotonicity of the matrix game, which holds regardless of how many equilibria
it has (the magnet's QRE path selects one). Dropping (I) weakens
Theorem~\ref{thm:main} from ``converges to $(w^\star,v^\star)$'' to ``converges
to a Nash of $A$'', which by interchangeability is the same statement about
$\varepsilon$.

\emph{It is needed for well-posedness in $s$.} The algorithm never faces $A$; it
faces the smoothed component matrix $\bar A = A + O(s^2)$ of \eqref{eq:Abar},
because a component is a Gaussian of width $s$ and not a point mass. If $A$'s
equilibrium is nondegenerate --- an open condition --- then $\bar A$'s is unique too
and lies $O(s^2)$ away, so the weights are continuous in the floor and the limit
can be named. If instead $A$ carries a continuum of equilibria, an arbitrarily
small perturbation can slide the solution to an extreme point of that continuum:
the weights become $O(1)$-sensitive to the smoothing scale and no statement of
the form ``the limit is near $w^\star$'' survives. \emph{(I) is a conditioning
requirement, not a convergence requirement.}

In the codebase's games it holds by construction rather than by luck: the
$K-1$ moment features of \texttt{MultiPointGame} make $w \mapsto \sum_j w_j
\varphi(p_j)$ injective on distributions over the peaks (a Vandermonde system in
the normalised peak coordinates, nonsingular because the peaks are distinct), so
exactly one weight vector matches the target moments --- see
Proposition~\ref{prop:moments}.
\end{remark}

\begin{remark}[(C) is much weaker than concavity]
Concavity of $(Q_\nu)_s$ on $\cA_j$ would be far too strong: a Gaussian bump
$h e^{-x^2/2w^2}$ smoothed to width $W = \sqrt{w^2+s^2}$ is concave only on
$|x| \le W$, yet mixture-MMD converges on the two-peak game from means far
outside that radius. Condition~\eqref{eq:star} only asks that the gradient have a
positive component \emph{toward} the branch point --- it permits arbitrarily flat
regions (small $\rho$), inflection, and anisotropy. What it forbids is exactly
(a) a second attractor inside the cell and (b) the gradient pointing away from
$\xi_j$ anywhere in the cell, which also makes $\cA_j$ forward invariant.
\end{remark}

\begin{remark}[(C) is uniform in the opponent, and that is the point]
Basins of a \emph{time-varying} flow are not forward invariant: a component in
the right basin at time $t$ can be in the wrong one at $t+1$ because the opponent
moved. So ``initialise each mean in the basin of its atom'' --- assumption A2 of
the previous theory --- is not a closable hypothesis on its own.
Condition~\eqref{eq:star} instead supplies a Lyapunov function
$V(a) = \tfrac12\|a - \xi_j\|^2$ that is \emph{the same function for every
opponent policy and every scale}, which is what makes the transport argument of
Section~\ref{sec:suff} go through.
\end{remark}

\begin{remark}[Flat landscapes are covered]\label{rem:flat}
If $u$ has no self-term ($u(a,b) = \ip{\varphi(a)}{C\psi(b)}$, the
``pure-coupling'' games of \texttt{gaussian\_magnet/THEORY.md}), then at
equilibrium $Q_{\nu^\star} \equiv \mathrm{const}$ and \eqref{eq:star} fails for
$\tau^{\mathrm g}_0 = 0$. With $\tau^{\mathrm g}_0 > 0$ the regularised landscape
\eqref{eq:reg} is strongly concave, \eqref{eq:star} holds with a single cell and
$\rho(r) = (\tau^{\mathrm g}_0/\bar\sigma^2) r$, and the class applies. The
conclusion then degrades in the expected way: the atoms are unidentified (any
moment-matched configuration is a fixed point), so convergence must be read as
convergence of the \emph{induced measure} and of exploitability, not of the
parameters. This is the correct unification of the ``magnet redundant'' (curved)
and ``magnet necessary'' (flat) regimes: they are the two extremes of a single
condition on the regularised landscape, and the generic game --- curvature in
some directions, flatness in others --- needs no separate treatment.
\end{remark}

\section{Why the class is the right one}

\subsection{Sharpness: (C) cannot be dropped}

\begin{proposition}[Spurious attractors are fatal]\label{prop:sharp}
Suppose there are $\hat a \in \cA$, $\hat b \in \cB$ and scales
$\hat s_0 \in \Sigma_0$, $\hat s_1 \in \Sigma_1$ such that, with
$\hat\mu = \cN(\hat a, \hat s_0^2)$ and $\hat\nu = \cN(\hat b, \hat s_1^2)$,
$\hat a$ is a nondegenerate local maximum of $(Q_{\hat\nu})_{\hat s_0}$ and
$\hat b$ a nondegenerate local maximum of $(W_{\hat\mu})_{\hat s_1}$. Let
$\hat z$ be the parameter configuration with all of player $0$'s means at
$\hat a$, all of player $1$'s at $\hat b$, and scales at $\hat s_i$. Then
$\hat z$ is a fixed point of \eqref{eq:cat}--\eqref{eq:gauss} with the magnet at
itself, and it is locally exponentially stable --- \emph{regardless of whether it
realises a Nash equilibrium}. Hence it has an open basin of initialisations from
which mixture-MMD does not converge to Nash.
\end{proposition}

\begin{proof}[Proof sketch]
Fixed point: all component $q$-values coincide (all components are identical), so
\eqref{eq:cat} is stationary; the mean gradients vanish because $\hat a$, $\hat b$
are critical points of the respective smoothed landscapes; the scale gradient is
the one-signed quantity of Lemma~\ref{lem:scale} and is inactive at the floor.
Stability: the own-parameter Hessian of each player is negative definite on the
active face (nondegenerate maximum), and by the block-skew structure of zero-sum
parametric games the symmetric part of the Jacobian of the pseudo-gradient field
is block diagonal with these two blocks; a negative definite symmetric part
remains Hurwitz after any positive-definite block-diagonal preconditioning,
because $G(G^{-1}J) + (G^{-1}J)^\top G = J + J^\top \prec 0$ exhibits
$V(z) = z^\top G z$ as a Lyapunov function. The magnet only adds
$-\eta\tau I$, improving the margin.
\end{proof}

Proposition~\ref{prop:sharp} is why the class must forbid \emph{all} spurious
local maxima of the effective landscape, not merely dominated ones. It also
corrects a reading of the previous theory: assumption A1 there (``every spurious
Gaussian bump is dominated by the atoms in both height and mass'') controls
which branch is the \emph{global} maximum at each scale, but the mean dynamics
perform local ascent and never compare global amplitudes. A dominated spurious
bump still captures any component that enters its basin. Height-and-mass
dominance is therefore neither necessary nor sufficient for convergence; it is a
proxy for branch persistence along the particular symmetric broad initialisation
used in those experiments.

\subsection{The three empirical failure edges are one condition}

The previous theory measured three independent-looking window edges. Under
Definition~\ref{def:class} they are three ways for \eqref{eq:star} to fail:

\begin{center}
\begin{tabular}{@{}l p{0.30\textwidth} p{0.40\textwidth}@{}}
\hline
edge & previous reading & failure of \eqref{eq:star}\\
\hline
$s_{\mathrm{dom}}$ (decoy) & spurious mass overtakes at large $s$ &
  $\xi_j(\nu,s)$ leaves $\cA_j$: wrong attractor\\
$s_{\mathrm{ej}}$ (ejection) & coupling expels a component &
  fails for some $\nu \in \cN_1$: opponent-dependent\\
$s_{\mathrm{reach}}$ (dead zone) & gradient $\sim e^{-d^2/2s^2}$ &
  \eqref{eq:star} \emph{holds}; $\rho$ is tiny --- a rate, not a failure\\
\hline
\end{tabular}
\end{center}

The demotion of $s_{\mathrm{reach}}$ from an assumption to a rate is substantive:
it says the dead zone never breaks convergence asymptotically, it only makes it
exponentially slow, which matches the observation that the boundary cases
converge at $150$k steps and not at $40$k. Correspondingly the class does not
assume a reach condition; a finite-horizon corollary would have to quantify
$\rho$.

\subsection{Sufficiency argument}\label{sec:suff}

We give the argument in the form we believe a proof should take: three
input-to-state-stable (ISS) blocks composed by a small-gain theorem. This is
\emph{not} a completed proof; the gaps are listed in Section~\ref{sec:gaps}.

\begin{lemma}[Mean block: tracking]\label{lem:mean}
Fix a component $k$ of player $0$ lying in cell $\cA_j$, and let
$r(t) = \|m_k(t) - \xi_j(t)\|$ where $\xi_j(t) = \xi_j(\nu_t, \sigma_k(t),
\bar a(t))$. Under (C), (W) and small enough $\eta_0$,
\begin{equation}\label{eq:iss}
  r(t+1) \;\le\; r(t) \;-\; \eta_0\, \underline\sigma^2 w_{\min}\, \rho\big(r(t)\big)
  \;+\; \big\|\xi_j(t+1) - \xi_j(t)\big\| \;+\; O(\eta_0^2).
\end{equation}
Consequently $\cA_j$ is forward invariant for $m_k$, and if the drift
$\|\dot\xi_j\| \to 0$ then $r(t) \to 0$.
\end{lemma}

\begin{proof}
Expand $\tfrac12 r^2$ along \eqref{eq:gauss} and apply \eqref{eq:star} to the
inner product; the drift term is the displacement of the target between steps.
Forward invariance follows because \eqref{eq:star} makes the radial component of
the step strictly inward, in particular on $\partial \cA_j$.
\end{proof}

\begin{lemma}[Scale block: scales fall]\label{lem:scale}
Let $\Lambda \in C^2(\R^n)$ and $\xi(s) \in \argmax_{\cA_j} \Lambda_s$ be
isolated. Then, with $\Lambda_s = \Lambda * \cN(0,s^2 I)$,
\begin{equation}
  \frac{d}{d(s^2/2)}\Big[\Lambda_s\big(\xi(s)\big)\Big]
  = \Delta \Lambda_s\big(\xi(s)\big) \;\le\; 0,
\end{equation}
with strict inequality at a nondegenerate maximum. Hence at (or near) a branch
point the log-scale natural gradient in \eqref{eq:gauss} is strictly negative and
$\sigma_k$ decreases monotonically to the floor $s^{\min}_0$.
\end{lemma}

\begin{proof}
$\partial_{(s^2/2)} \Lambda_s = \Delta \Lambda_s$ is the heat equation; the
envelope theorem kills the $\dot\xi$ term; $\Delta \Lambda_s \le 0$ at an interior
maximum.
\end{proof}

Lemma~\ref{lem:scale} replaces the empirical ``self-annealing'' story with a
proof, and it holds in any dimension. It also explains why an explicit annealing
schedule is never needed for $\sigma \to s^{\min}$: a schedule only changes
\emph{which} branch you are attracted to on the way down, which is precisely the
mechanism by which annealing delivers a component into a decoy in games outside
the class.

\begin{lemma}[Weight block]\label{lem:weight}
Let $\bar A(t)$ be the component matrix \eqref{eq:Abar}. The categorical updates
\eqref{eq:cat} are tabular MMD on $\bar A(t)$; for fixed $\bar A$ and fixed
magnet they contract to the $\tau^{\mathrm{cat}}$-regularised equilibrium (QRE) at
rate $1 - \Theta(\eta\tau^{\mathrm{cat}})$, and the magnet refresh drives the QRE
to the Nash of $\bar A$. For time-varying $\bar A(t)$ the block is ISS with gain
$O(1/\tau^{\mathrm{cat}})$ in $\|\dot{\bar A}\|$ (the QRE map is
$1/\tau^{\mathrm{cat}}$-Lipschitz in the payoff matrix).
\end{lemma}

\begin{theorem}[Sufficiency, modulo the gaps of §\ref{sec:gaps}]\label{thm:main}
Let $u \in \mathfrak{G}(M,\Sigma,\cN,\rho)$ in the sense of
Definition~\ref{def:class}. Then the iterates of
\eqref{eq:cat}--\eqref{eq:gauss} satisfy: every component mean converges to a
branch point of its cell; every scale converges to the floor; the induced
component matrix converges to the smoothed atom matrix
$\bar A_{jl} \to A^{s}_{jl}$; the categorical heads converge to the unique Nash
of that matrix; and the induced policy pair is an $\varepsilon$-Nash equilibrium
of the continuous game with, \emph{under} (B),
\begin{equation}\label{eq:bias}
  \varepsilon = O\big((s^{\min}_0)^2 + (s^{\min}_1)^2\big),
\end{equation}
the constant being $\tfrac12 \sum_j w^\star_j\, |\nabla^2 Q_{\nu^\star}(p_j)|$
for player $0$ and analogously for player $1$.
\end{theorem}

\begin{proof}[Proof skeleton]
\emph{Step 1 (transport).} By (K) each cell contains a mean; by
Lemma~\ref{lem:mean} each mean stays in its cell and tracks its branch with error
governed by the branch drift. By (D) the round-trip gain between the two players'
branch maps is $g_0 g_1 < 1$, so the small-gain theorem applied to the two mean
blocks gives that the joint tracking error contracts up to the input from the
weight and scale blocks.

\emph{Step 2 (scales).} By Lemma~\ref{lem:scale} the scales are eventually
monotonically decreasing once the tracking error is small enough that the
components sit near their branch maxima; they converge to the floor, and their
drift $\to 0$, removing one input to Step 1.

\emph{Step 3 (weights).} By Lemma~\ref{lem:weight} the categorical heads track
the Nash of the drifting component matrix with gain $O(\|\dot{\bar A}\|/
\tau^{\mathrm{cat}})$. By Steps 1--2, $\dot{\bar A} \to 0$, so the weights
converge; by (M) the limiting matrix is the atom matrix $A^s$. Nondegeneracy is
an \emph{open} condition, so the $O(s^2)$-perturbed matrix $A^s$ still has a
unique fully-mixed equilibrium and, by (I), the weights converge to
$(w^\star, v^\star) + O(s^2)$. Without (I) they still converge, to \emph{some}
Nash of $A^s$ (Remark~\ref{rem:nondeg}).

\emph{Step 4 (bias).} At the limit each component sits at an atom with scale
$s^{\min}$, so the exploitability is the smoothing deficit
$\max_a Q_{\nu} - \E_{\pi} Q_\nu$ per player. Under (B) a second-order Taylor
expansion of $Q_{\nu^\star}$ at $p_j$ against the Gaussian gives
$\tfrac12 \sum_j w_j |\nabla^2 Q(p_j)| (s^{\min})^2 + O(s^4)$, i.e.\
\eqref{eq:bias}. In the Gaussian-bump case this is
$\sum_j w_j h_j\big(1 - w_j/\sqrt{w_j^2 + s^2}\big)$ exactly.

\emph{Step 5 (duplicates).} If a cell receives several components, all of them
converge to the same branch point; the induced matrix has duplicated rows, whose
Nash equilibrium is non-unique in $w$ but unique as a measure on the atoms. The
conclusion is stated for the induced measure and is unaffected.
\end{proof}

\begin{corollary}[Overcapacity helps, inside the class]
Under Definition~\ref{def:class}, taking $K_i \gg M_i$ with a space-filling
initialisation makes the covering condition in (K) automatic, at no cost to the
conclusion (Step 5). This is the opposite of the behaviour observed on the decoy
game, where extra components are captured by the spurious attractor of
Proposition~\ref{prop:sharp} --- a game \emph{outside} the class. Overcapacity is
therefore the cheap way to discharge the only genuinely initialisation-dependent
hypothesis, provided the game is in the class.
\end{corollary}

\subsection{The floor, and how exactness is recovered}\label{sec:floor}

The floor is a hypothesis, not an approximation, and it is worth being explicit
about what it buys and what it costs. It buys three things:

\begin{enumerate}
\item \textbf{Breadth in the payoff}: Remark~\ref{rem:reg} --- bounded measurable
  instead of $C^2$.
\item \textbf{Breadth in the landscape}: condition (C) is imposed from
  $s^{\min}$ upward, not from $0$. In one dimension Gaussian smoothing never
  creates an extremum, so a landscape carrying structure finer than the floor ---
  ripple, discretisation noise, a Monte-Carlo estimate of $D$ --- can have very
  many local maxima at scale $0$ and exactly $M_0$ at scale $s^{\min}$. Such a
  game is inside the class with a floor and outside it without. \emph{A positive
  floor makes the class strictly broader}, which is the main reason to adopt it.
\item \textbf{Uniform constants}, hence a checkable condition (a compact scale
  range) and a linear local rate $1 - \Theta(\eta\tau)$ rather than a degrading
  one.
\end{enumerate}

It costs the resolution requirement of Lemma~\ref{lem:resolution}: structure
finer than the floor is invisible, so atoms finer than the floor are
unrepresentable.

The apparent cost --- that the conclusion is $\varepsilon$-Nash rather than Nash
--- is not real, because the floor is not what drives the bias down. The scale
dynamics do that by themselves, and releasing the floor turns
Theorem~\ref{thm:main}'s $O(s^2)$ into a rate in $t$.

\begin{corollary}[Unfloored: $\varepsilon = O(1/t)$]\label{cor:unfloored}
Suppose the hypotheses of Theorem~\ref{thm:main} have delivered transport and
equilibration --- every mean at its branch point, weights at the matrix-game
equilibrium --- and the floor is then released ($s^{\min} \to 0$). Assume (B).
Writing $\kappa_j = |\Delta Q_{\nu^\star}(p_j)| > 0$, the log-scale update
\eqref{eq:gauss} gives, from $\partial_{\log \sigma}(Q)_\sigma(\xi) = \sigma^2
\Delta Q(\xi) + O(\sigma^4)$,
\begin{equation}
  \dot\sigma_k = -c_k\,\sigma_k^3 + O(\sigma^5),
  \qquad c_k = \tfrac{\eta}{2}\, w_k\, \kappa_k ,
\end{equation}
hence $\sigma_k(t) = \big(\sigma_k(0)^{-2} + 2c_k t\big)^{-1/2} = \Theta(t^{-1/2})$
and, by Step 4 of Theorem~\ref{thm:main},
\begin{equation}
  \varepsilon(t) \;=\; \Theta(1/t).
\end{equation}
Transport survives the collapse, but only marginally: the mean contraction rate
carries the preconditioner $\sigma^2 \sim 1/t$, so the accumulated contraction is
$\int^t \sigma^2 \,d\tau = \Theta(\log t)$ --- divergent, hence still convergent,
but logarithmically.
\end{corollary}

\begin{lemma}[The Gaussian magnet time-rescales the Gaussian-head flow]
\label{lem:brake}
Let the Gaussian-head magnet have temperature $\tau^{\mathrm g}$ and be refreshed
every $T$ steps, and set
\begin{equation}\label{eq:brake}
  \beta \;\coloneqq\; \frac{1 - e^{-\eta \tau^{\mathrm g} T}}
                            {\eta \tau^{\mathrm g} T} \;\in\; (0, 1] ,
  \qquad \beta \to 1 \text{ as } \tau^{\mathrm g} \to 0 .
\end{equation}
Then, to leading order in the per-cycle displacement, \emph{both} Gaussian-head
coordinates are slowed by the same factor $\beta$:
\begin{enumerate}
\item[(i)] the cycle-averaged mean velocity is $\beta$ times its
  $\tau^{\mathrm g} = 0$ value;
\item[(ii)] Corollary~\ref{cor:unfloored} holds verbatim with
  $c_{\mathrm{eff}} = \beta\, c$.
\end{enumerate}
Consequently the $(m,\sigma)$ flow with magnet is, to this order, the
magnet-free flow under the time reparameterisation $t \mapsto \beta t$.
\end{lemma}

\begin{proof}
(i) Within a cycle write $x = m - \bar m$. The magnet contributes
$-\eta\tau^{\mathrm g}(\sigma^2/\bar\sigma^2)\,x \approx
-\eta\tau^{\mathrm g} x$ to the update, so $\dot x = g - \eta\tau^{\mathrm g}x$
with $g$ the magnet-free step; from $x(0) = 0$,
$x(T) = (g/\eta\tau^{\mathrm g})(1 - e^{-\eta\tau^{\mathrm g}T})$, and the
average velocity $x(T)/T$ is $\beta g$.

(ii) Within one cycle $\sigma$ is constant to $O(1/t)$. Writing
$x = \rho - \bar\rho$, the magnet's log-scale gradient is
$B(1 - e^{2x})$ with $B = \tfrac{\eta}{2}\tau^{\mathrm g}$ (from
$\partial_\rho \KL = -1 + \sigma^2/\bar\sigma^2$), so with $g = c\,\sigma^2$,
$\dot x = -g + B(1 - e^{2x}) = A - Be^{2x}$, $A = B - g$. In $u = e^{2x}$ this
is logistic, $\dot u = 2u(A - Bu)$, and from $u(0) = 1$,
$x(T) = \tfrac12 \log\!\big[(1-\varepsilon)/(1 - \varepsilon e^{-2AT})\big]$ with
$\varepsilon = g/B$, giving decay per cycle
$|x(T)| = \tfrac{\varepsilon}{2}(1 - e^{-2BT}) + O(\varepsilon^2)$, i.e.\
$c_{\mathrm{eff}} = \beta c$ on dividing by $T$ (note
$2BT = \eta\tau^{\mathrm g}T$, the same exponent as in (i) --- which is why the
two coordinates share one factor).
\end{proof}

The brake \emph{saturates} inside the refresh cycle whenever
$\eta\tau^{\mathrm g}T \gtrsim 1$, which is the regime the defaults sit in
($\eta\tau T = 2$). A linear-response estimate --- which is what one writes down
first --- misses this and gives $\beta = 1/2$, an error of $15\%$ at the defaults
and unbounded as $\eta\tau^{\mathrm g}T$ grows.

\begin{remark}[The bias exponent is a certificate of (B)]\label{rem:biasexp}
Without (B) the exponent degrades rather than the conclusion failing. At a corner
maximum ($Q \sim -|a - p|$ near the atom) the deficit is
$\E|\cN(0,s^2)| = s\sqrt{2/\pi} = O(s)$, so $\varepsilon = O(s^{\min})$ and, in
Corollary~\ref{cor:unfloored}, $\varepsilon = O(t^{-1/2})$. So the exponent $2$
measured empirically (log--log slope $2.002$ in the previous theory's
experiment~5) is evidence that the landscape is $C^2$ at its atoms, not a
universal law.
\end{remark}

\subsection{Numerical verification}\label{sec:verify}

Corollary~\ref{cor:unfloored} and Lemma~\ref{lem:brake} are sharp enough to test
directly, and \texttt{theory2/check\_scale\_law.py} does so on
\texttt{TWO\_POINT} ($h = 1$, $\bar w = 0.1$, coupling $1$, $K = 2$), whose
predicted free constant is $c = \tfrac{\eta}{2} w_k h/\bar w^2 = 1.25$. The run
loop is a transcription of \texttt{idealized\_mmd.run} exposing two knobs the
original hard-codes (the floor, and a separate $\tau^{\mathrm g}$); run~A
verifies it reproduces \texttt{run} bit-for-bit and reproduces the previous
theory's experiment~2 to all printed digits. Exploitability is measured on a
$3\times10^5$-point grid, because at the repo's default $4001$ points neither
peak is a grid point and the best-response value is short of the true maximum by
$\sim 7\times10^{-6}$ per player --- the same order as the quantities being
measured.

\begin{center}
\begin{tabular}{lccc}
\hline
run ($2\times 10^6$ steps, floor $10^{-12}$) & fitted $c$ & predicted $c$ &
  ratio\\
\hline
B: $\tau^{\mathrm g} = 0$ (free decay) & $1.249835$ & $1.25$ & $0.9999$\\
C: $\tau^{\mathrm g} = 0.2$, atom init & $0.541092$ & $0.540415$ & $1.0013$\\
D: $\tau^{\mathrm g} = 0.2$, spread init $(-0.3, 1.3)$ & $0.541092$ &
  $0.540415$ & $1.0013$\\
\hline
\end{tabular}
\end{center}

\noindent Fitting $\sigma^{-2}$ against $t$ over the last two decades gives
$R^2 \ge 1 - 3\times10^{-8}$ in every run: the trajectory is affine in
$\sigma^{-2}$, exactly as $\dot\sigma = -c\sigma^3$ requires. The log--log slopes
are $\sigma$: $-0.502$ to $-0.506$ (predicted $-0.5$) and $\varepsilon$:
$-1.003$ to $-1.011$ (predicted $-1$), and $\varepsilon = 100\,\sigma^2$ holds to
five significant figures throughout, confirming that the residual is pure
representation bias. Run~D differs from run~C only in the fitted intercept: a
transport phase changes where the law starts, not the law.

Lemma~\ref{lem:brake}'s factor was checked by sweeping $\tau^{\mathrm g}$ and $T$
over a $16\times$ range in $\eta\tau^{\mathrm g}T$; the two configurations sharing
$\eta\tau^{\mathrm g}T = 0.5$ (namely $\tau^{\mathrm g}=0.2, T=50$ and
$\tau^{\mathrm g}=0.05, T=200$) collapse onto the same $c$, confirming the factor
depends only on that product:

\begin{center}
\begin{tabular}{cccccc}
\hline
$\tau^{\mathrm g}$ & $T$ & $\eta\tau^{\mathrm g}T$ & factor & $c$ measured &
  ratio\\
\hline
$0.2$ & $50$ & $0.5$ & $0.78694$ & $0.98715$ & $1.0035$\\
$0.05$ & $200$ & $0.5$ & $0.78694$ & $0.98429$ & $1.0006$\\
$0.2$ & $200$ & $2.0$ & $0.43233$ & $0.54092$ & $1.0009$\\
$0.2$ & $800$ & $8.0$ & $0.12496$ & $0.15583$ & $0.9976$\\
$0.8$ & $200$ & $8.0$ & $0.12496$ & $0.15586$ & $0.9979$\\
$0$ & $200$ & $0$ & $1$ & $1.24967$ & $0.9997$\\
\hline
\end{tabular}
\end{center}

\begin{remark}[The floor is typically inactive]\label{rem:inactive}
The same runs settle a reading of the previous theory's measurements. Its
experiment~2 reports, at $2\times10^4$ steps with the floor at $10^{-3}$, a
converged $\sigma = 0.0069105$ --- which is $6.9\times$ the floor. \emph{The floor
was never active}: that run's scale was still falling, at the rate
\eqref{eq:brake} predicts, and its reported exploitability $0.0047574$ is the
smoothing deficit $2h(1 - \bar w/\sqrt{\bar w^2 + \sigma^2}) = 0.0047$ at that
$\sigma$. So the residual exploitabilities quoted throughout that work are
snapshots of a $\Theta(1/t)$ decay rather than floor-set constants, and comparing
two runs' residuals compares their step counts unless the floor is verified
active in both.
\end{remark}

\subsection{The Gaussian magnet and the class}\label{sec:magnet}

Because condition (C) is imposed on the regularised landscape \eqref{eq:reg},
$\tau^{\mathrm g}$ sits inside the class definition and it is fair to ask whether
the class is larger with it. The answer is directional, and it is different in
the two regimes of the previous theory's magnet dichotomy.

\paragraph{Flat directions: strictly broader.} For a pure-coupling game
$u(a,b) = \ip{\varphi(a)}{C\psi(b)}$ with no self-term, at equilibrium
$Q_{\nu^\star} \equiv \mathrm{const}$, so $\nabla Q \equiv 0$ and \eqref{eq:star}
fails at every $a$ for $\tau^{\mathrm g} = 0$: these games are \emph{outside} the
class. With $\tau^{\mathrm g} > 0$ the quadratic in \eqref{eq:reg} makes
$Q^{\mathrm{reg}}$ strongly concave, (C) holds with a single cell and
$\rho(r) = (\tau^{\mathrm g}/\bar\sigma^2)\,r$, and they are admitted. So the
magnet does enlarge the class --- by exactly the games whose landscape supplies no
curvature. The price is the one recorded in Remark~\ref{rem:flat}: the branch
point is $\bar a$, the previous iterate, so (M) has no content, the atoms are
unidentified, and the conclusion is convergence of the induced measure rather
than of the parameters.

\paragraph{Curved directions: no broadening, and a cost.} Here the magnet cannot
buy membership, because the obstruction to (C) is a spurious local maximum and
the magnet does not remove one:

\begin{proposition}[Spurious attractors are magnet-invariant]\label{prop:magnetinv}
The configuration $\hat z$ of Proposition~\ref{prop:sharp} is a fixed point of
\eqref{eq:cat}--\eqref{eq:gauss} for \emph{every} $\tau^{\mathrm g} \ge 0$, and
its local contraction rate is \emph{increasing} in $\tau^{\mathrm g}$.
\end{proposition}

\begin{proof}
At $\hat z$ the magnet is refreshed to $\hat z$ itself, so $\bar a = \hat a$ and
the magnet's mean gradient $-\tau^{\mathrm g}(\hat a - \bar a)/\bar\sigma^2$
vanishes identically; the fixed-point property is therefore unchanged. The
magnet adds $-\eta\tau^{\mathrm g}I$ to the Jacobian, which strictly improves the
stability margin of an already-stable fixed point.
\end{proof}

So the magnet makes a trap \emph{deeper}, never shallower --- the analytic form
of the previous theory's ``no magnet fixes transport'', and of its observation
that a stronger magnet is worse on the multidimensional decoy game.

\paragraph{What the cost is, exactly.} By Lemma~\ref{lem:brake} the magnet
time-rescales the whole Gaussian-head flow by $\beta$. Two consequences for the
hypotheses:
\begin{itemize}
\item \emph{Asymptotic basins are unchanged.} A time reparameterisation does not
  alter which initialisations converge, so the covering condition (K), read
  asymptotically, is insensitive to $\tau^{\mathrm g}$.
\item \emph{Finite-horizon reach shrinks by $\beta$.} At a fixed step budget, the
  set of initialisations that actually arrive contracts, and by
  Corollary~\ref{cor:unfloored} the bias at a fixed step count is worse by
  $1/\beta$. This is the sense --- the only sense --- in which the magnet narrows
  things.
\end{itemize}

Run~F of \texttt{check\_scale\_law.py} measures the transport time from the
spread initialisation $(-0.3, 1.3)$ against the prediction $1/\beta$:

\begin{center}
\begin{tabular}{cccc}
\hline
$\tau^{\mathrm g}$ & $\beta$ & measured $t_{\mathrm{arrive}}$ ratio &
  predicted $1/\beta$\\
\hline
$0$ & $1$ & $1.00$ & $1.00$\\
$0.2$ & $0.43233$ & $2.67$ & $2.31$\\
$2.0$ & $0.05000$ & $20.18$ & $20.00$\\
\hline
\end{tabular}
\end{center}

\noindent The $\tau^{\mathrm g} = 2$ row is the clean test ($0.9\%$); the
$\tau^{\mathrm g} = 0.2$ row is looser because early transport happens at large
$\sigma$, before the per-cycle displacement is small enough for the leading-order
rescaling to apply.

\begin{remark}[The dichotomy is per direction, not per game]\label{rem:perdir}
At an interior equilibrium the own-parameter Hessian is negative semidefinite, so
each own-direction is either curved (magnet redundant, and costly) or flat
(magnet necessary). The generic game mixes the two --- e.g.\ a two-dimensional
action with a well in one coordinate and matching pennies in the other --- and the
correct prescription, ``magnet on the flat null-space only'', is not expressible
with a scalar $\tau^{\mathrm g}$. If any direction is flat, one pays the friction
$\beta$ in \emph{all} of them. A direction-aware magnet (a $\KL$ weighted by a
projector onto the null-space of the own-Hessian) is the obvious remedy and, as
far as we know, untried.
\end{remark}

\subsection{What is not proved}\label{sec:gaps}

\begin{enumerate}
\item \textbf{The small-gain composition.} Lemmas~\ref{lem:mean}--\ref{lem:weight}
  are individually routine; composing three ISS blocks requires the gains to be
  simultaneously small, and the weight block's gain $O(1/\tau^{\mathrm{cat}})$
  pushes toward large $\tau^{\mathrm{cat}}$ while the bias argument pushes toward
  the magnet refresh driving regularisation away. The ordering of limits
  (refresh period $T$ versus gain margins) is not settled.
\item \textbf{Quantitative $\rho$.} The class gives asymptotic convergence. A
  finite-horizon statement --- the honest form of the old $s_{\mathrm{reach}}$
  edge --- needs $\rho$ explicitly, and for Gaussian-bump landscapes $\rho$ is
  exponentially small at long range.
\item \textbf{Verification of (C).} Checking a condition quantified over all
  opponent policies and all scales is the practical cost of the class;
  Section~\ref{sec:check} reduces it in useful special cases but does not
  eliminate it.
\item \textbf{The flat case} (Remark~\ref{rem:flat}) is covered by the definition
  but Theorem~\ref{thm:main}'s Steps 1 and 3 must be reread in moment
  coordinates, since the fixed-point set is a manifold rather than a point.
\item \textbf{Steps 1 and 2 are not sequential.} Theorem~\ref{thm:main} is
  written as though transport completes and then the scales fall, but the mean
  update carries the preconditioner $\sigma^2$, so the scale collapse of
  Lemma~\ref{lem:scale} \emph{throttles} transport as it proceeds: at
  $\sigma = 0.1$ the mean contraction is $\approx 0.025$ per step, at
  $\sigma = 10^{-3}$ it is $\approx 2.5 \times 10^{-6}$ --- four orders of
  magnitude, i.e.\ a component that has not arrived before the scale collapses is
  effectively frozen where it stands. With a floor the rate is merely small and
  the argument is unaffected; unfloored (Corollary~\ref{cor:unfloored}) both the
  contraction and the branch drift vanish together and the ISS estimate
  \eqref{eq:iss} needs their \emph{ratio} controlled, not each separately. This
  is the one place where the two-block decomposition is genuinely coupled and I
  have not verified the ratio is favourable in general.
\end{enumerate}

\section{Checkable sufficient conditions}\label{sec:check}

Condition (C) quantifies over $\nu \in \cN_1$. Two standard structures make that
quantifier finite-dimensional and the whole condition checkable.

\subsection{Finite-rank coupling}

Suppose
\begin{equation}\label{eq:G1}
  u(a,b) = D_0(a) - D_1(b) + \ip{\varphi(a)}{C\,\psi(b)},
  \qquad \varphi : \cA \to \R^{m_0},\ \psi : \cB \to \R^{m_1},
\end{equation}
with $D_i, \varphi, \psi \in C^2$. Then $Q_\nu(a) = D_0(a) + \ip{\varphi(a)}{Cy}
+ \mathrm{const}$ with $y = \E_\nu \psi$, so the opponent enters through
$m_1$ numbers and $\cN_1$ collapses to the compact convex \emph{moment body}
$Y = \conv \psi(\cB) \subset \R^{m_1}$. Condition (C) becomes: for all $y \in Y$
and all $s \in \Sigma_0^{n_0}$, the function
$D_{0,s} + \ip{\varphi_s(\cdot)}{Cy}$ is star-concave toward $\xi_j(y,s)$ on
$\cA_j$. Two consequences.

\begin{proposition}[Curvature-dominates-coupling]\label{prop:cap}
Suppose $D_{0,s}$ satisfies \eqref{eq:star} on $\cA_j$ with modulus $\rho_0$ and
that the branch $\xi_j^{0}(s)$ of $D_{0,s}$ has curvature $\nabla^2 D_{0,s} \preceq
-\alpha_0 I$ near it. If
\begin{equation}
  \sup_{y \in Y} \|C y\| \cdot \sup_{\cA_j} \|\nabla \varphi_s\| \;<\;
  \inf_{r>0} \rho_0(r),
\end{equation}
then \eqref{eq:star} holds for the coupled landscape with modulus
$\rho_0 - \|Cy\|\|\nabla\varphi_s\|$, and the branch drift satisfies
$\|\partial \xi_j/\partial y\| \le \|\nabla \varphi_s\|\,\|C\|/\alpha_0$ by the
implicit function theorem. Combining the two players, the small-gain condition
\eqref{eq:smallgain} is implied by
\begin{equation}\label{eq:cap}
  \alpha_0\, \alpha_1 \;>\; \|C\|^2\, \sup\|\nabla\varphi\|\, \sup\|\nabla\psi\| .
\end{equation}
\end{proposition}

For the codebase games, $u = D(a) - D(b) + c\ip{f(a)}{f(b)}$ with Gaussian bumps
of height $h$ and width $w$, so $\alpha \approx h/w^2$ and \eqref{eq:cap} reads
$c\,\|f'\|^2 < (h/w^2)^2$ --- the same $O(h/w^2)$ coupling cap that the previous
theory measured locally, now doing global work.

\begin{example}[The cap on the codebase games, with numbers]\label{ex:cap}
\texttt{TWO\_POINT} has $h = 1$, $w = 0.1$, hence $\alpha = h/w^2 = 100$; its
feature is $f(a) = 2a - 1$, so $\|f'\| = 2$; and $c = 1$. Then
\[
  g_0 \;=\; \frac{c\,\|f'\|}{\alpha} \;=\; \frac{2}{100} \;=\; 0.02,
  \qquad
  g_0 g_1 \;=\; 4 \times 10^{-4} \;\lll\; 1 ,
\]
and \eqref{eq:cap} would first bite at $c \approx \alpha/\|f'\| = 50$. The
coupling in that game is therefore three orders of magnitude away from being able
to destabilise the branch chase --- which is why no experiment at $c = 1$ has ever
shown coupling-induced trouble, and why the two players' targets settle rather
than chase.

Consistency check against the previous theory's coupling sweep, which used the
three-peak game ($w = 0.08$, so $\alpha \approx 156$) with a quadratic feature:
\eqref{eq:cap} predicts trouble from $c \approx \alpha/\|f'\| \approx 78$, and the
certified region there was measured flat in $c$ through $c = 8$ and degrading
between $c = 32$ and $c = 300$. Same scale and same $O(h/w^2)$ form, from a
condition that --- unlike the measured one --- is a statement about the whole
transit rather than a neighbourhood of the equilibrium.
\end{example}

\begin{proposition}[Nondegeneracy from moments]\label{prop:moments}
In the symmetric-coupling case $C = cI$, $\psi = \varphi$, condition (I) holds
iff the vectors $\{\varphi(p_1),\dots,\varphi(p_{M_0})\}$ are affinely
independent and
$0 \in \mathrm{relint}\,\conv\{\varphi(p_1),\dots,\varphi(p_{M_0})\}$. This is the
general form of the Vandermonde argument in \texttt{MultiPointGame}: with
$\varphi_j = \iota^j - t_j$ the moment map is a Vandermonde system in the
normalised peak coordinates, injective because the peaks are distinct.
\end{proposition}

\subsection{Gaussian-bump landscapes in one dimension}

If $D_0 = \sum_b h_b \exp(-(x-c_b)^2/2w_b^2)$ then:
\begin{itemize}
\item \emph{Mode counting suffices at one scale --- the floor.} In one dimension
  Gaussian smoothing never creates a local extremum (scale-space causality), so
  the mode count of $D_{0,s}$ is non-increasing in $s$. Condition (C)'s
  ``exactly $M_0$ attractors at every scale in $\Sigma_0$'' therefore reduces to
  a count at the \emph{bottom} of the range plus a merge check: \emph{$D_{0,s}$
  has exactly $M_0$ local maxima at $s = s^{\min}_0$, and no two of them merge
  below $s^{\max}_0$}. Both are one-dimensional root counts.
\item \emph{Structure finer than the floor is free.} Because the count is taken
  at $s^{\min}_0$ and not at $0$, a landscape of the form ``$M_0$ clean bumps
  $+$ ripple of wavelength $\ll s^{\min}_0$'' passes: the ripple contributes
  arbitrarily many local maxima to $D_0$ and none to $D_{0,s^{\min}}$. This is
  the concrete content of the claim that the floor broadens the class, and it is
  what makes the class usable when $D$ is only available as an estimate.
\item \emph{The decoy game is excluded immediately}: its landscape has three
  local maxima (two peaks and the decoy) while $|\supp\mu^\star| = 2$, so (M)
  fails at $M_0 = 3 \ne 2$ --- no smoothing or dominance computation needed. This
  is the cleanest available statement of why that game is a counterexample.
\item \emph{The two- and three-peak games are included} (number of maxima
  $=$ support size), subject to \eqref{eq:cap} and (I).
\end{itemize}

\subsection{Multidimensional actions}\label{sec:multidim}

Definition~\ref{def:class} is dimension-free: cells are subsets of $\R^{n}$,
\eqref{eq:star} is an inner-product inequality, the scale range is a box
$\Sigma^n$ for the diagonal-covariance policy (and \eqref{eq:star} is imposed for
every anisotropic scale in that box), and Lemma~\ref{lem:scale} uses the
Laplacian. What changes is verification and cost:

\begin{enumerate}
\item \textbf{No scale-space causality.} For $n \ge 2$, Gaussian smoothing
  \emph{can} create local maxima, and a sum of $\ge 3$ Gaussians can have more
  modes than components. So the mode count at $s^{\min}$ does not certify the
  ladder; (C) must be imposed on the entire range, and the honest certificate is
  a numerical mode-ladder check, not a bump-wise inequality.
\item \textbf{Wide-low structures are exponentially worse.} An isotropic bump
  $(h,w)$ smoothed at scale $s$ has amplitude $h\,(w^2/(w^2+s^2))^{n/2}$, so its
  effective mass is $h w^n$. Any spurious structure that is wide relative to the
  atoms gains an exponential-in-$n$ advantage at large scale, which in the
  language of the class means the branch maps $\xi_j(\nu,s)$ leave their cells at
  a scale $s^{\max}$ decreasing in $n$: the admissible scale range shrinks with
  dimension.
\item \textbf{Covering costs more.} (K) asks that every cell contain an initial
  mean; in $\R^n$ this needs $K$ growing with the number of cells and with the
  geometry, which is the multidimensional price of the only
  initialisation-dependent hypothesis.
\item \textbf{Bias grows linearly in $n$}: the trace in Step 4 of
  Theorem~\ref{thm:main} runs over $n$ axes, so $\varepsilon = O(n\,s^2)$ and the
  floor must shrink like $1/\sqrt n$ for a fixed target.
\item \textbf{Separable games and the policy class.} If
  $u(a,b) = \sum_{d=1}^{n} u_d(a_d, b_d)$, then (C) factorises: verifying the
  one-dimensional condition per axis suffices, and the equilibrium is a product
  of the per-axis equilibria. But then the \emph{policy} must factorise too: a
  joint mixture with a shared categorical head makes the per-axis component
  values add, so one axis's weight preference drags the weights of all axes
  simultaneously, and (I) --- nondegeneracy of the induced matrix game's equilibrium
  --- is stated for a joint matrix of size $K^n$ that a $K$-component joint
  mixture cannot represent. The correct statement is a dichotomy:
  \begin{itemize}
  \item genuinely $M$-atom equilibrium in $\R^n$ $\Rightarrow$ joint
    $K \ge M$-component mixture, class applies verbatim;
  \item separable game with per-axis support $M$ $\Rightarrow$ \emph{per-axis
    factorised} categorical head with $K \ge M$ per axis, class applies per axis.
    A joint mixture with a shared head is outside the analysis, and is the
    documented source of the dimension-amplified weight-starvation failure.
  \end{itemize}
\end{enumerate}

\section{Relation to the previous theory}

\begin{center}
\begin{tabular}{p{0.30\textwidth} p{0.62\textwidth}}
\hline
previous object & status here\\
\hline
Lifted skewness (Lemma 1) & used, unchanged; it is why the categorical block
  \eqref{eq:cat} is globally convergent.\\
Block-skew structure (Lemma 2) & used only \emph{locally}, in
  Proposition~\ref{prop:sharp}, and repaired for the Fisher metric via the
  Lyapunov identity $GM + M^\top G = J + J^\top$. Global monotonicity of the
  parametric field is not claimed: the flow that is run is $G^{-1}F$, and
  monotonicity of $F$ does not survive a state-dependent preconditioner.\\
Local island theorem (Thm.\ 4) & the $r \to 0$ endpoint of
  Lemma~\ref{lem:mean}; the linear rate $1 - \Theta(\eta\tau)$ is
  Lemma~\ref{lem:weight}.\\
$O(s^2)$ bias (Lemma 3) & Step 4 of Theorem~\ref{thm:main}, restated for a
  general landscape via $\nabla^2 Q$, and \emph{conditional on} (B): the
  exponent $2$ certifies $C^2$ at the atoms rather than holding universally
  (Remark~\ref{rem:biasexp}). Released from the floor it becomes
  $\varepsilon = \Theta(1/t)$ (Corollary~\ref{cor:unfloored}), a law verified numerically to four digits in
\S\ref{sec:verify}.\\
Variance floor $s^{\min}$ & promoted from an implementation detail to a
  hypothesis: it is what allows $u$ merely bounded measurable
  (Remark~\ref{rem:reg}) and what makes (C) tolerant of landscape structure
  finer than the floor (§\ref{sec:floor}). Its only cost is the new
  Lemma~\ref{lem:resolution}.\\
A1 scale dominance & \emph{demoted}: neither necessary nor sufficient
  (§\ref{sec:suff}); replaced by (C)+(M).\\
A2 reach & \emph{demoted} to a rate statement carried by $\rho$.\\
A3 coupling caps & Proposition~\ref{prop:cap}, now serving the global argument.\\
A4 step size & (S), unchanged.\\
A5 capacity/genericity & split into (K) capacity+covering, (I) nondegeneracy,
  (W) weight non-degeneracy --- with (W) identified as an artifact of the
  preconditioner (Remark~\ref{rem:wk}).\\
$\sigma$-niceness & (C) is its uniform, opponent-aware, dimension-free
  strengthening: the ``path property'' (no branch jumping) is automatic because a
  star-concave cell has one attractor, and ``local strong concavity'' is
  weakened to star-concavity with modulus $\rho$.\\
Magnet dichotomy (curvature vs.\ magnet) & absorbed: (C) is imposed on the
  \emph{regularised} landscape \eqref{eq:reg}, so curved and flat directions are
  handled by one condition, per direction, with no case split.\\
\hline
\end{tabular}
\end{center}

\section{Summary}

The class is: \emph{zero-sum, bounded measurable, compact; the family of smoothed
magnet-regularised effective landscapes has exactly $M_i$ one-point-attracting
cells, uniformly over opponent policies and over the scale range
$[s^{\min}_i, s^{\max}_i]$ with $s^{\min}_i > 0$, with branch maps whose
round-trip drift gain is a contraction; the atoms are resolved at the floor; the
branch endpoints at equilibrium and at the floor scale are the Nash atoms; the
induced atom matrix game has a unique equilibrium; capacity $K_i \ge M_i$ with an
initialisation meeting every cell; and the usual step-size and magnet
conditions.} Nothing is assumed about symmetry between the players, about the
algebraic form of the payoff, about smoothness of the payoff, about concavity, or
about the action dimension.

The strict variance floor is load-bearing in both directions. It is what lets the
payoff be merely bounded measurable and what lets the landscape carry structure
finer than the floor, so it \emph{widens} the class rather than restricting it to
the implemented algorithm; and it costs only Lemma~\ref{lem:resolution}, the requirement that
the equilibrium atoms be resolved at that scale. Exactness is not sacrificed:
released from the floor, the scale dynamics are $\sigma \sim t^{-1/2}$ and the
representation bias vanishes as $\Theta(1/t)$
(Corollary~\ref{cor:unfloored}), a law verified numerically to four digits in
\S\ref{sec:verify}.

The reason to believe it suffices is Lemma~\ref{lem:decomp}: the categorical
heads solve a finite matrix game exactly and globally, so the only thing left to
assume is that the means find the atoms --- and \eqref{eq:star} is the weakest
standard condition under which a preconditioned gradient method finds a
prescribed point, made uniform in the opponent so that it survives the game
dynamics. The reason to believe it cannot be much broader is
Proposition~\ref{prop:sharp}: the qualitative failure of \eqref{eq:star} produces
an exponentially stable non-Nash configuration with an open basin, so no
initialisation-robust class can permit it.

\end{document}
